%%%%%%%%%%%%%%%%%%%%%%%%%%%%%%%%%%%%%%%%%%%%%%%%%%%%%%%%%%%%%%%%%%%%%%
%%%%%%%%% Select one of the options, and comment the rest of them

%%%%%%%%%% Option 1:  to compile with pdflatex : parameter "t" - to align to the top
\documentclass[professionalfont]{beamer}
%sans font?

%%%%%%%%%% Option 3: to create handout for print
%\documentclass[t,handout]{beamer}
%\usepackage{pgfpages}              % to put several slides on one page
%\pgfpagesuselayout{2 on 1}[a4paper, border shrink=5mm]             % 2 slides on 1 page
%\pgfpagesuselayout{4 on 1}[a4paper,landscape, border shrink=5mm]   % 4 slides on 1 page, and landscaped
\usepackage{xcolor}
\usepackage{tikz}

%%%%%%%%%%%%%%%%%%%%%%%%%%%%%%%%%%%%%%%%%%%%%%%%%%%%%%%%%%%%%%%%%%%%
%%%%%%%%%%%%%% Select the Theme %%%%%%%%%%%%%%%%%%%%%%%%%%%%%%%%%%%
\usetheme{Dresden}     % OK
%\usetheme{Berlin}
%\usetheme{Bergen}      % NO
%\usetheme{Boadilla}    % NO
%\usetheme{Copenhagen}  % NO
%\usetheme{Hannover}    % NO
%\usetheme{Luebeck}     % NO
%\usetheme{Marburg}     % NO
%\usetheme{Pittsburgh}  % NO
%\usetheme{default}
%\usetheme{Singapore}   % OK
%\usetheme{boxes}
%\usecolortheme{structure}
%\usecolortheme{rose}
%\usecolortheme{beaver}


\definecolor{mymaroon}{cmyk}{0.0, 1.0, 1.0, 0.498}
\definecolor{myblue}{cmyk}{1.0, 1, 0, 0.5}
\definecolor{mygreen}{cmyk}{100, 0, 100, 50}
\setbeamercolor*{palette secondary}{use=structure,fg=white,bg=myblue}
\setbeamercolor*{palette tertiary}{use=structure,fg=white,bg=mymaroon}

%\usepackage{beamerthemesplit}              %
\beamertemplateballitem % fancy bullets and numbering

\setbeamertemplate{navigation symbols}{}   % suppress navigation symbols
\addtobeamertemplate{frametitle}{}{%
	\begin{tikzpicture}[remember picture,overlay]
	\node[anchor=center, yshift=-15pt, xshift=-5pt] at (current page.north) 
	{\includegraphics[height=1.1cm]{../images/ANL_black}\hspace{10cm}};
	
	\node[anchor=north east, yshift=1pt, xshift=-5pt] at (current page.north east) 
	{\includegraphics[height=0.8cm]{../images/IIT_Logo_blk}};
	\end{tikzpicture}}


% other possibilities to include LOGO. it puts it in RLC

%
%\pgfdeclareimage[width=1cm]{logo}{../images/IIT_Logo}
%\logo{\pgfuseimage{logo}}


% load additional packages
\usepackage{graphicx}
\usepackage{amsmath}
\usepackage{amssymb}
\usepackage{amsthm}
\usepackage{graphicx}
\usepackage{url}
\usepackage{color}


% Your Abbreviations
\newcommand\bE{{\mathbb{E}}}
\newcommand\bR{{\mathbb{R}}}
\newcommand\bH{{\mathbf{H}}}
% End abbreviations


%%%%%%%%%%%%%%%%%%%%% to edit the main text below
%NOTES ON SOME TECHNICS
%%%% Box %%%%%%%%%%%%%%%%%%%%%%%%%%%%%%%%%%%%%%%%%%%%%%%
%{\fbox{ \parbox[t]{10cm}{ SOME TEXT }}}

%%% include a picture. The file should be with extention EPS, e.g. FILENAME.EPS
%\begin{figure}[h]
%\centering
%\includegraphics[width=.7\linewidth]{FILENAME}
%\caption{{\footnotesize PUT_CAPTION }}
%\end{figure}

%\subtitle{}
%\institute[ANL/IIT]{Argonne National Laboratory\\Illinois Institute of Technology}

\title[Space Charge 2017]{Study of space charge dominated beams at the AWA rf photoinjector}
\author[N.Neveu]{{\Large Nicole Neveu}}
\institute[ANL, IIT] % (optional, but mostly needed)
{   Illinois Institute of Technology \\
	Argonne National Laboratory \\
    \url{nneveu@anl.gov} 
}
% - Use the \inst command only if there are several affiliations.
% - Keep it simple, no one is interested in your street address.
\date{ \today \\
\includegraphics[width=3cm,keepaspectratio]{../images/Argonne_cmyk_black}%
\hfill \hfill \hfill%
\includegraphics[width=4cm,keepaspectratio]{../images/IIT_logo}%
}

%\date[IIT, April 2009]{
%           Space Charge 2017 \\ Oc 18, 2009  }



\begin{document}


\begin{frame}
  \titlepage
\end{frame}


%\begin{frame}{Outline of the talk}
%  \tableofcontents
%  % You might wish to add the option [pausesections]
%\end{frame}


% Structuring a talk is a difficult task and the following structure
% may not be suitable. Here are some rules that apply for this
% solution:

% - Exactly two or three sections (other than the summary).
% - At *most* three subsections per section.
% - Talk about 30s to 2min per frame. So there should be between about
%   15 and 30 frames, all told.

% - A conference audience is likely to know very little of what you
%   are going to talk about. So *simplify*!
% - In a 20min talk, getting the main ideas across is hard
%   enough. Leave out details, even if it means being less precise than
%   you think necessary.
% - If you omit details that are vital to the proof/implementation,
%   just say so once. Everybody will be happy with that.

\section{Starting Point}

\subsection{Install Beamer}

\begin{frame}{Beamer Installation}

\begin{itemize}

\item the official web site \\[.2cm]
\quad \url{http://latex-beamer.sourceforge.net} \\[.4cm]

\item download the latest release and find User Guide \\[.4cm]

\item User Guide, files this talk and some additional materials \\[.2cm]
\quad \url{ math.iit.edu/~igor/beamer} \\[.4cm]

\item Mik\TeX\&WinEdt will recognize (most likely) the missing packages and install them automatically
when you try to compile a presentation \\[.4cm]

\item try to compile  the SimpleFile.tex (\LaTeX \  $\Rightarrow$ DVIPS $\Rightarrow$ Open PS).
It should compile and
you should be able to see the slides. If not, work on installation.
\end{itemize}
\end{frame}


\begin{frame}{How to compile. Ya!}

\begin{itemize}
\item \LaTeX \  $\Rightarrow$ DVIPS $\Rightarrow$ Open PS \\[.5cm]
\item PDF\LaTeX \ $\Rightarrow$ view PDF \\[.5cm]
\item Also after PS you can get the PDF: ps2pdf
\end{itemize}

\end{frame}

\subsection{Simple Examples}
\begin{frame}[containsverbatim, shrink]{Example}
\begin{verbatim}
\documentclass{beamer}
\title{Example Presentation Created with the Beamer Package}
\author{Your Name Here} \date{\today}

\begin{document}
\frame{\titlepage}

\section{Introduction}
\subsection{Overview of the Beamer Class}

\begin{frame}{Features of the Beamer Class}
\begin{itemize}
  \item<1-> Normal LaTeX class.
  \item<2-> Easy overlays.
  \item<3-> No external programs needed.
  \end{itemize}
\end{frame}
\end{document}

\end{verbatim}
\end{frame}


\begin{frame}[containsverbatim, shrink]{DocumentClass}

General Form: \begin{verbatim}\documentclass[options]{beamer}

%%%%%%%%%% Option 1:  to compile with pdflatex
\documentclass[mathserif]{beamer}

%%%%%%%%%% Option 2:  to compile latex, dvips, ps2pdf
\documentclass[t,dvips, mathserif]{beamer}

%%%%%%%%%% Option 3: to create handouts for print
\documentclass[t,handout]{beamer}
\usepackage{pgfpages} % to put several slides on one page
\pgfpagesuselayout{2 on 1}[a4paper, border shrink=5mm]
% 2 slides on 1 page
%\pgfpagesuselayout{4 on 1}[a4paper,landscape, border shrink=5mm]
% 4 slides on 1 page, and landscaped
\end{verbatim}
\end{frame}



\begin{frame}[containsverbatim]{Select the Theme/ColorTheme}
\begin{verbatim}
\usetheme{Dresden}      %\usetheme{Berlin}
%\usetheme{Luebeck}     %\usetheme{Marburg}
%\usetheme{Pittsburgh}  %\usetheme{default}

%\usecolortheme{structure} %\usecolortheme{rose}
\usecolortheme{beaver}
\end{verbatim}

\begin{itemize}
\item And many more. Read the manual.
\item {\bf Rule of Thumbs: Simple is better.}
\item A math presentation should be a math presentation and
not a slide-show for Vouge or Glamour Magazine.
\end{itemize}
\end{frame}



\begin{frame}[containsverbatim]{Cover Slide}

\begin{verbatim}
....
\title{A Good Title}
\author{That's your name here}
\institute{You work somewhere, I guess}

\begin{document}

% make first slide/frame, cover slide
\begin{frame}
    \titlepage
\end{frame}

\end{document}

\end{verbatim}

\end{frame}



\begin{frame}[containsverbatim]

Make fancy numbering and bullets:
\begin{verbatim}
    \beamertemplateballitem
\end{verbatim}

\smallskip

Nice splitting of the sections and subsections \begin{verbatim}
    \usepackage{beamerthemesplit}
\end{verbatim}

\begin{block}{Summary}
Put something in a block
\end{block}

\begin{verbatim}
    \begin{block}{Summary}
    Put something in a block
    \end{block}
\end{verbatim}

\end{frame}

\subsection{blocks}

\begin{frame}[containsverbatim]{Blocks, Theorems, Examples}

\begin{beamerboxesrounded}[shadow=true]{Theorem}
Prove that $x^n+y^n=z^n$.
\end{beamerboxesrounded}

\begin{example}
For example $3^2+4^2=5^2$
\end{example}

\begin{verbatim}
    \begin{beamerboxesrounded}[shadow=true]{Theorem}
    Prove that $x^n+y^n=z^n$.
    \end{beamerboxesrounded}

    \begin{example}
    For example $3^2+4^2=5^2$
    \end{example}
\end{verbatim}
\end{frame}


\subsection{overlays}


\begin{frame}{Transition}

Pause command
\begin{itemize}
\pause \item One Step
\pause \item Two
\pause \item The end.
\end{itemize}



Or More complicated (not necessarily better)\pause

\begin{enumerate}
\item<5-6> Beatles
\item<6-7>  Pick Floyd
\item<-7> Led Zeppelin
\item<-8>  Lenon
\end{enumerate}
\end{frame}

\begin{frame}[containsverbatim]{The code}
\begin{verbatim}
Pause command
\begin{itemize}
\pause \item One Step
\pause \item Two
\pause \item The end.
\end{itemize}

Or More complicated (not necessarily better)\pause

\begin{enumerate}
\item<5-6> Beatles
\item<6-7>  Pick Floyd
\item<-7> Led Zeppelin
\item<-8>  Lenon
\end{enumerate}

\end{verbatim}
\end{frame}


\begin{frame}{Finally}
\begin{block}{Remark}
Avoid complicated overlays and too much animation. Usually ``pause'' is more than enough.
\end{block}
\bigskip
Other overlay commands: ONLY, UNCOVER, INVISIBLE, ALT,

\pause
\begin{center}
\bigskip
{\Large  The End \\[.5cm] Good Luck with your slides!}
\end{center}
\end{frame}






\end{document}